% !TeX root = ../apunte_mecanica_clasica.tex
\IfFileExists{../apunte_mecanica_clasica.tex}{%
  \documentclass[../apunte_mecanica_clasica.tex]{subfiles}%
}{%
  \documentclass[apunte_mecanica_clasica.tex]{subfiles}%
}
\begin{document}

\chapter{Operadores Diferenciales Clásicos}

\section{Operador Gradiante}


\subsection{Forma cartesiana y notación covariante}

El operador gradiente es bien conocido en la física-matemática.  
En coordenadas cartesianas se escribe como
\begin{equation*}
    \vec{\nabla} = \hat{\imath}\,\frac{\partial}{\partial x}
    + \hat{\jmath}\,\frac{\partial}{\partial y}
    + \hat{k}\,\frac{\partial}{\partial z} \,.
\end{equation*}

Al actuar sobre una función arbitraria $f(x,y,z)$ se obtiene:
\begin{equation*}
    \vec{\nabla} f =
    \hat{\imath}\,\frac{\partial f}{\partial x}
    + \hat{\jmath}\,\frac{\partial f}{\partial y}
    + \hat{k}\,\frac{\partial f}{\partial z} \,.
\end{equation*}

La idea principal es escribir el operador gradiente en forma 
\textit{covariante}, esto es, que posea un índice arriba (contravariante) 
y uno abajo (covariante).  

Entonces escribimos el operador gradiente como:
\begin{equation*}
    \vec{\nabla} = \hat{\imath}\,\frac{\partial}{\partial x}
    + \hat{\jmath}\,\frac{\partial}{\partial y}
    + \hat{k}\,\frac{\partial}{\partial z} \,,
\end{equation*}
lo cual equivale a
\begin{equation*}
    \vec{\nabla} = \hat{e}_1 \frac{\partial}{\partial x^1}
    + \hat{e}_2 \frac{\partial}{\partial x^2}
    + \hat{e}_3 \frac{\partial}{\partial x^3} \,,
\end{equation*}
o bien, en notación compacta,
\begin{equation*}
    \vec{\nabla} = \sum_{i=1}^3 \hat{e}_i \,\frac{\partial}{\partial x^i} \,.
\end{equation*}

Aquí, tanto $\hat{e}_i$ como $\dfrac{\partial}{\partial x^i}$ llevan 
el índice en posición inferior.

Usando la identidad $\hat{e}_i = \hat{E}^i$, obtenida para coordenadas
cartesianas, se tiene que
\begin{equation*}
    \vec{\nabla} = \sum_{i=1}^3 \hat{E}^i \,\frac{\partial}{\partial x^i} \,.
\end{equation*}


\newpage

\subsection{Definición general}

\begin{apuntebloque}{Definición}
Se define el \textbf{Operador Gradiente} en el espacio vectorial 
$\mathbb{R}^3$ como
\begin{equation*}
    \vec{\nabla} = \sum_{i=1}^3 \vec{E}^i \, \frac{\partial}{\partial x^i} \,.
\end{equation*}
\end{apuntebloque}

Si usamos la fórmula general para las bases duales $\{\vec{E}^i\}_{i=1}^3$ se obtiene:

\begin{equation*}
    \vec{E}^i = \sum_{j=1}^3 g^{ij} \, \vec{e}_j \,.
\end{equation*}

Por lo tanto, el operador gradiente se escribe como:
\begin{equation*}
    \vec{\nabla} = \sum_{i=1}^3 \sum_{j=1}^3 g^{ij} \, \vec{e}_j \, \frac{\partial}{\partial x^i} \,.
\end{equation*}


\subsection{Ejemplos}

\begin{apuntebloque}{Ejemplo: Operador Gradiente}
\textbf{a) Coordenadas Cartesianas:} En coordenadas cartesianas el operador gradiente es bien conocido:
\begin{equation*}
    \vec{\nabla} f = \hat{\imath}\,\frac{\partial f}{\partial x}
                   + \hat{\jmath}\,\frac{\partial f}{\partial y}
                   + \hat{k}\,\frac{\partial f}{\partial z} \,.
\end{equation*}

\textbf{b) Coordenadas Esféricas:} En coordenadas esféricas el operador gradiente se escribe como:
\begin{align*}
    \vec{\nabla} f &= \sum_{i=1}^3 \vec{E}^i \frac{\partial f}{\partial x^i} \\[4pt]
    &= \vec{E}^1 \frac{\partial f}{\partial x^1}
     + \vec{E}^2 \frac{\partial f}{\partial x^2}
     + \vec{E}^3 \frac{\partial f}{\partial x^3} \\[4pt]
    &= \vec{e}_r \,\frac{\partial f}{\partial r}
     + \frac{\vec{e}_\phi}{r^2 \sin^2\theta}\,\frac{\partial f}{\partial \phi}
     + \frac{\vec{e}_\theta}{r^2}\,\frac{\partial f}{\partial \theta} \,.
\end{align*}

Usando los vectores unitarios asociados a las bases coordenadas esféricas:
\begin{equation*}
    \vec{e}_r = \hat{r}, \qquad
    \vec{e}_\phi = r \sin\theta \,\hat{\phi}, \qquad
    \vec{e}_\theta = r\,\hat{\theta},
\end{equation*}
se obtiene finalmente:
\begin{equation*}
    \vec{\nabla} f = \hat{r}\,\frac{\partial f}{\partial r}
                   + \frac{\hat{\phi}}{r \sin\theta}\,\frac{\partial f}{\partial \phi}
                   + \frac{\hat{\theta}}{r}\,\frac{\partial f}{\partial \theta} \,.
\end{equation*}
\end{apuntebloque}

\section{Operador Divergencia}


\subsection{Derivación en coordenadas generales}

En coordenadas cartesianas:

\begin{align*}
    \vec{\nabla}\cdot\vec{A} 
    &= \left( \hat{\imath}\,\frac{\partial}{\partial x} 
           + \hat{\jmath}\,\frac{\partial}{\partial y} 
           + \hat{k}\,\frac{\partial}{\partial z}\right)
       \cdot \left(A_x \hat{\imath} + A_y \hat{\jmath} + A_z \hat{k}\right) \\[6pt]
    &= \frac{\partial A_x}{\partial x} 
     + \frac{\partial A_y}{\partial y} 
     + \frac{\partial A_z}{\partial z}.
\end{align*}

En un sistema de coordenadas general, con

\begin{equation*}
    \vec{A} = \sum_{j=1}^3 A^j \vec{e}_j.
\end{equation*}

\begin{align*}  
    \vec{\nabla}\cdot\vec{A} 
    &= \sum_{i=1}^3 \vec{E}^i \frac{\partial}{\partial x^i} 
       \cdot \left( \sum_{j=1}^3 A^j \vec{e}_j \right) \\[6pt]
    &= \sum_{i=1}^3 \sum_{j=1}^3 
       \vec{E}^i \cdot \left( 
            \frac{\partial A^j}{\partial x^i} \vec{e}_j 
            + A^j \frac{\partial \vec{e}_j}{\partial x^i}
       \right) \\[6pt]
    &= \sum_{i=1}^3 \sum_{j=1}^3 
       \frac{\partial A^j}{\partial x^i} \, \vec{E}^i \cdot \vec{e}_j
       + \sum_{i=1}^3 \sum_{j=1}^3 A^j 
          \vec{E}^i \cdot \left( \sum_{k=1}^3 \Gamma^k{}_{ji} \vec{e}_k \right) \\[6pt]
    &= \sum_{i=1}^3 \sum_{j=1}^3 
       \frac{\partial A^j}{\partial x^i} \,\delta^i{}_j
       + \sum_{i=1}^3 \sum_{j=1}^3 A^j \sum_{k=1}^3 \Gamma^k{}_{ji}\,\delta^i{}_k \\[6pt]
    &= \sum_{i=1}^3 \sum_{j=1}^3 
       \frac{\partial A^j}{\partial x^i} \,\delta^i{}_j
       + \sum_{i=1}^3 \sum_{j=1}^3 A^j \,\Gamma^i{}_{ji}.
\end{align*}

\begin{align*}
    \Rightarrow \quad 
    \vec{\nabla}\cdot\vec{A}
    &= \sum_{i=1}^3 \frac{\partial A^i}{\partial x^i}
       + \sum_{j=1}^3 A^j \sum_{i=1}^3 \Gamma^i{}_{ji}.
\end{align*}

Usando la identidad
\begin{equation*}
    \sum_{i=1}^3 \Gamma^i{}_{ji} 
    = \frac{\partial}{\partial x^j} \ln \sqrt{g}, 
    \qquad g = \det(g_{ij}),
\end{equation*}
se obtiene:
\begin{align*}
    \vec{\nabla}\cdot\vec{A} 
    &= \sum_{i=1}^3 \frac{\partial A^i}{\partial x^i}
       + \sum_{j=1}^3 A^j \frac{\partial}{\partial x^j}\ln\sqrt{g} \\[6pt]
    &= \sum_{i=1}^3 \frac{\partial A^i}{\partial x^i}
       + \sum_{i=1}^3 A^i \frac{\partial}{\partial x^i}\ln\sqrt{g} \\[6pt]
    &= \sum_{i=1}^3 \frac{\partial A^i}{\partial x^i}
       + \sum_{i=1}^3 A^i \frac{1}{\sqrt{g}}
          \frac{\partial \sqrt{g}}{\partial x^i} \\[6pt]
    &= \frac{1}{\sqrt{g}} \sum_{i=1}^3 
       \left( \sqrt{g}\,\frac{\partial A^i}{\partial x^i}
            + A^i \frac{\partial \sqrt{g}}{\partial x^i} \right) \\[6pt]
    &= \frac{1}{\sqrt{g}} \sum_{i=1}^3 
       \frac{\partial}{\partial x^i}\!\left(\sqrt{g}\,A^i\right).
\end{align*}


\subsection{Componentes en bases unitarias}

En este punto se vuelve importante usar las bases coordenadas 
\textbf{unitarias}, puesto que se desea que las componentes $A^i$ 
tengan la misma dimensionalidad que el vector $\vec{A}$.  

Por ejemplo, en coordenadas esféricas el vector $\vec{A}$ se escribe como:
\begin{equation*}
    \vec{A} = A_r \,\vec{e}_r + A_\phi \,\vec{e}_\phi + A_\theta \,\vec{e}_\theta,
\end{equation*}
donde los vectores coordenados $\vec{e}_\phi$ y $\vec{e}_\theta$ poseen la 
dimensionalidad del radio:
\begin{equation*}
    [\vec{e}_r] = [\vec{e}_\phi] = [r] = \text{Longitud}.
\end{equation*}

Por lo tanto, las componentes $A_\phi$ y $A_\theta$ no poseen la misma 
dimensionalidad que el vector $\vec{A}$.

Sin embargo, los vectores bases unitarios 
\begin{equation*}
    \hat{e}_i = \frac{\vec{e}_i}{\|\vec{e}_i\|},
\end{equation*}
son adimensionales.  

Entonces escribimos el vector $\vec{A}$ en términos de los vectores bases, tal que
\begin{equation*}
    \vec{A} = \sum_{i=1}^3 A^i \,\vec{e}_i 
            = \sum_{i=1}^3 A^i \,\|\vec{e}_i\| \,\hat{e}_i.
\end{equation*}

Por lo tanto,
\begin{equation*}
    A^i = \frac{\tilde{A}^i}{\|\vec{e}_i\|}, 
    \qquad \vec{A} = \sum_{i=1}^3 \tilde{A}^i \,\hat{e}_i,
\end{equation*}
donde $\tilde{A}^i$ son las componentes asociadas a los vectores unitarios $\hat{e}_i$.

Por lo tanto, la divergencia del vector $\vec{A}$ se escribirá como
\begin{equation*}
    \vec{\nabla}\cdot \vec{A} 
    = \frac{1}{\sqrt{g}} \sum_{i=1}^3 
      \frac{\partial}{\partial x^i} \left( \sqrt{g}\, \frac{\tilde{A}^i}{\|\vec{e}_i\|} \right),
\end{equation*}
donde $A^i$ son las componentes del vector respecto a las bases coordenadas
y $\|\vec{e}_i\|$ son las normas de los vectores base.

\begin{apuntebloque}{Ejemplo: Operador Divergencia}

\textbf{a) Coordenadas Cartesianas:}  
El determinante de $\delta_{ij}$ es $1$, luego
\begin{align*}
    \vec{\nabla}\cdot \vec{A} 
    &= \frac{1}{\sqrt{1}} \sum_{i=1}^3 
       \frac{\partial}{\partial x^i} \left( \sqrt{1}\, A^i \right) \\[6pt]
    &= \sum_{i=1}^3 \frac{\partial A^i}{\partial x^i} \\[6pt]
    &= \frac{\partial A_x}{\partial x}
     + \frac{\partial A_y}{\partial y}
     + \frac{\partial A_z}{\partial z}.
\end{align*}

\vspace{0.5cm}

\textbf{b) Coordenadas Esféricas:}  
El determinante de la métrica asociada a las coordenadas esféricas es
\begin{equation*}
    g = r^4 \sin^2 \theta,
\end{equation*}
    
y las normas de sus vectores bases coordenados:
\[
    \|\vec{e}_r\| = 1, \qquad 
    \|\vec{e}_\phi\| = r \sin \theta, \qquad 
    \|\vec{e}_\theta\| = r,
\]

entonces:
\begin{align*}
    \vec{\nabla}\cdot \vec{A} 
    &= \frac{1}{\sqrt{r^4 \sin^2 \theta}} 
       \sum_{i=1}^3 \frac{\partial}{\partial x^i} 
       \left( \sqrt{r^4 \sin^2 \theta}\; 
              \frac{\tilde{A}^i}{\|\vec{e}_i\|} \right) \\[6pt]
    &= \frac{1}{r^2 \sin \theta} 
       \sum_{i=1}^3 \frac{\partial}{\partial x^i} 
       \left( r^2 \sin \theta \; \frac{\tilde{A}^i}{\|\vec{e}_i\|} \right)\\[6pt]
     &= \frac{1}{r^2 \sin \theta} 
        \frac{\partial}{\partial x^1} 
       \left( r^2 \sin \theta \; \frac{\tilde{A}^1}{\|\vec{e}_1\|} \right)+\frac{1}{r^2 \sin \theta} 
        \frac{\partial}{\partial x^2} 
       \left( r^2 \sin \theta \; \frac{\tilde{A}^2}{\|\vec{e}_2\|} \right)\\[6pt]
     &\quad + \frac{1}{r^2 \sin \theta} 
        \frac{\partial}{\partial x^3} 
       \left( r^2 \sin \theta \; \frac{\tilde{A}^3}{\|\vec{e}_3\|} \right)\\[8pt]  
     &= \frac{1}{r^2 \sin \theta} 
        \frac{\partial}{\partial r} 
       \left( r^2 \sin \theta \; \frac{\tilde{A}^r}{\|\vec{e}_r\|} \right)+\frac{1}{r^2 \sin \theta} 
        \frac{\partial}{\partial \phi} 
       \left( r^2 \sin \theta \; \frac{\tilde{A}^{\phi}}{\|\vec{e}_{\phi}\|} \right)\\[6pt]
     &\quad + \frac{1}{r^2 \sin \theta} 
        \frac{\partial}{\partial \theta} 
       \left( r^2 \sin \theta \; \frac{\tilde{A}^{\theta}}{\|\vec{e}_{\theta}\|} \right)\\[8pt]
         &= \frac{1}{r^2 \sin \theta} 
        \frac{\partial}{\partial r} 
       \left( r^2 \sin \theta \; \frac{\tilde{A}^r}{1} \right)+\frac{1}{r^2 \sin \theta} 
        \frac{\partial}{\partial \phi} 
       \left( r^2 \sin \theta \; \frac{\tilde{A}^{\phi}}{r \sin\theta} \right)\\[6pt]
     &\quad + \frac{1}{r^2 \sin \theta} 
        \frac{\partial}{\partial \theta} 
       \left( r^2 \sin \theta \; \frac{\tilde{A}^{\theta}}{r} \right)\\[8pt]  
    &= \frac{1}{r^2 } 
       \frac{\partial}{\partial r}\left(r^2\,\tilde{A}^r \right) 
     + \frac{1}{r \sin \theta} 
       \frac{\partial}{\partial \phi}\left(\,\tilde{A}^\phi \right)+ \frac{1}{r \sin \theta} 
       \frac{\partial}{\partial \theta}\left(\sin \theta \,\tilde{A}^\theta \right).
\end{align*}

Por lo tanto:
\begin{equation*}
    \vec{\nabla}\cdot \vec{A} 
    = \frac{1}{r^2}\frac{\partial}{\partial r}\!\left(r^2 \tilde{A}^r\right)
    + \frac{1}{r \sin \theta}\frac{\partial}{\partial \phi}\!\left(\tilde{A}^\phi\right)
    + \frac{1}{r \sin \theta}\frac{\partial}{\partial \theta}\!\left(\sin \theta\, \tilde{A}^\theta\right).
\end{equation*}

\end{apuntebloque}

\section{Operador Laplaciano}


\subsection{Definición y forma general}

\begin{apuntebloque}{Definición}
Se define el \textbf{operador Laplaciano} como
\begin{equation*}
    \nabla^2 f = \vec{\nabla}\cdot \vec{\nabla} f,
\end{equation*}
donde $f$ es una función arbitraria que depende de las coordenadas.
\end{apuntebloque}

Para obtener la forma general en coordenadas curvilíneas, definimos
\begin{align*}
    \vec{A} = \vec{\nabla} f
    &= \sum_{i=1}^3 \vec{E}^i \,\frac{\partial f}{\partial x^i} \\[6pt]
    &= \sum_{i=1}^3 \sum_{j=1}^3 g^{ij}\,\vec{e}_j \,\frac{\partial f}{\partial x^i}.
\end{align*}

Reordenando los índices:
\begin{align*}
    \vec{A} 
    &= \sum_{j=1}^3 \sum_{i=1}^3 g^{ji}\,\vec{e}_i \,\frac{\partial f}{\partial x^j} \\[6pt]
    &= \sum_{i=1}^3 \sum_{j=1}^3 g^{ij} \frac{\partial f}{\partial x^j}\,\vec{e}_i \\[6pt]
    &= \sum_{i=1}^3 A^i \,\vec{e}_i,
\end{align*}
donde
\begin{equation*}
    A^i = \sum_{j=1}^3 g^{ij}\,\frac{\partial f}{\partial x^j}.
\end{equation*}

Entonces,
\begin{align*}
    \nabla^2 f 
    &= \vec{\nabla}\cdot \vec{\nabla} f \\[6pt]
    &= \vec{\nabla}\cdot \vec{A} \\[6pt]
    &= \frac{1}{\sqrt{g}} \sum_{i=1}^3 \frac{\partial}{\partial x^i} 
       \left(\sqrt{g}\, A^i\right) \\[6pt]
    &= \frac{1}{\sqrt{g}} \sum_{i=1}^3 \frac{\partial}{\partial x^i} 
       \left(\sqrt{g}\, \sum_{j=1}^3 g^{ij}\,\frac{\partial f}{\partial x^j}\right).
\end{align*}

\[
    \therefore \quad 
    \nabla^2 f = \frac{1}{\sqrt{g}} \sum_{i=1}^3 \sum_{j=1}^3 
    \frac{\partial}{\partial x^i}\!\left(\sqrt{g}\, g^{ij}\,\frac{\partial f}{\partial x^j}\right).
\]


\subsection{Ejemplo en coordenadas esféricas}

\begin{apuntebloque}{Ejemplo: Operador Laplaciano}
\textbf{a) Coordenadas esféricas:}  
Tenemos que
\begin{align*}
    \nabla^2 f 
    &= \frac{1}{\sqrt{r^4 \sin^2\theta}} 
       \sum_{i=1}^3 \sum_{j=1}^3 
       \frac{\partial}{\partial x^i} 
       \left( r^4 \sin^2\theta \, g^{ij} \,\frac{\partial f}{\partial x^j} \right) \\[6pt]
    &= \frac{1}{r^2 \sin\theta} 
       \sum_{i=1}^3 \sum_{j=1}^3 
       \frac{\partial}{\partial x^i} 
       \left( r^2 \sin\theta \, g^{ij} \,\frac{\partial f}{\partial x^j} \right).
\end{align*}

Evaluando explícitamente:
\begin{align*}
    \nabla^2 f 
    &= \frac{1}{r^2 \sin\theta} 
       \left[
          \frac{\partial}{\partial r} \Big( r^2 \sin\theta \, g^{11}\,\frac{\partial f}{\partial r} \Big)
          + \frac{\partial}{\partial \phi} \Big( r^2 \sin\theta \, g^{22}\,\frac{\partial f}{\partial \phi} \Big) \right. \\[6pt]
    &\qquad \left. + \frac{\partial}{\partial \theta} \Big( r^2 \sin\theta \, g^{33}\,\frac{\partial f}{\partial \theta} \Big)
       \right].
\end{align*}

Recordemos que en coordenadas esféricas:
\begin{equation*}
    g^{11} = 1, 
    \qquad g^{22} = \frac{1}{r^2 \sin^2\theta}, 
    \qquad g^{33} = \frac{1}{r^2}.
\end{equation*}
    

Por lo tanto,
\begin{align*}
    \nabla^2 f 
    &= \frac{1}{r^2 \sin\theta} 
       \left[
          \frac{\partial}{\partial r}\!\Big(r^2 \sin\theta \,\frac{\partial f}{\partial r}\Big)
          + \frac{\partial}{\partial \phi}\!\Bigg(r^2 \sin\theta \,\frac{1}{r^2 \sin^2\theta}\,\frac{\partial f}{\partial \phi}\Bigg) \right. \\[6pt]
    &\qquad \left. + \frac{\partial}{\partial \theta}\!\Bigg(r^2 \sin\theta \,\frac{1}{r^2}\,\frac{\partial f}{\partial \theta}\Bigg)
       \right].
\end{align*}

Simplificando cada término:
\begin{align*}
    \nabla^2 f
    &= \frac{1}{r^2} \frac{\partial}{\partial r}\!\Big(r^2 \frac{\partial f}{\partial r}\Big)
       + \frac{1}{r^2 \sin^2\theta} \frac{\partial^2 f}{\partial \phi^2}
       + \frac{1}{r^2 \sin\theta} \frac{\partial}{\partial \theta}\!\Big(\sin\theta \frac{\partial f}{\partial \theta}\Big).
\end{align*}
\end{apuntebloque}

\newpage

\end{document}
